\subsubsection{The weighted version}

So far we have just been counting paths; but we can easily introduce weights
to obtain a more general result:

\begin{theorem}
[LGV lemma, lattice weight version]\label{thm.lgv.kpaths.wt}Let $K$ be a
commutative ring.

For each arc $a$ of the digraph $\mathbb{Z}^{2}$, let $w\left(  a\right)  $ be
an element of $K$. We call this element $w\left(  a\right)  $ the
\emph{weight} of $a$.

For each path $p$ of $\mathbb{Z}^{2}$, define the \emph{weight} $w\left(
p\right)  $ of $p$ by%
\[
w\left(  p\right)  :=\prod_{a\text{ is an arc of }p}w\left(  a\right)  .
\]


For each path tuple $\mathbf{p}=\left(  p_{1},p_{2},\ldots,p_{k}\right)  $,
define the \emph{weight} $w\left(  \mathbf{p}\right)  $ of $\mathbf{p}$ by%
\[
w\left(  \mathbf{p}\right)  :=w\left(  p_{1}\right)  w\left(  p_{2}\right)
\cdots w\left(  p_{k}\right)  .
\]


Let $k\in\mathbb{N}$. Let $\mathbf{A}=\left(  A_{1},A_{2},\ldots,A_{k}\right)
$ and $\mathbf{B}=\left(  B_{1},B_{2},\ldots,B_{k}\right)  $ be two
$k$-vertices. Then,%
\[
\det\left(  \left(  \sum_{p:A_{i}\rightarrow B_{j}}w\left(  p\right)  \right)
_{1\leq i\leq k,\ 1\leq j\leq k}\right)  =\sum_{\sigma\in S_{k}}\left(
-1\right)  ^{\sigma}\sum_{\substack{\mathbf{p}\text{ is a nipat}\\\text{from
}\mathbf{A}\text{ to }\sigma\left(  \mathbf{B}\right)  }}w\left(
\mathbf{p}\right)  .
\]
Here, \textquotedblleft$p:A_{i}\rightarrow B_{j}$\textquotedblright\ means
\textquotedblleft$p$ is a path from $A_{i}$ to $B_{j}$\textquotedblright.
\end{theorem}

Clearly, Proposition \ref{prop.lgv.kpaths.count} is the particular case of
Theorem \ref{thm.lgv.kpaths.wt} when $K=\mathbb{Z}$ and $w\left(  a\right)
=1$ for all arcs $a$ (because in this case, all the weights $w\left(
p\right)  $ and $w\left(  \mathbf{p}\right)  $ of paths and path tuples equal
$1$, and therefore the sums over paths or nipats become the \#s of paths or nipats).

\begin{proof}
[Proof of Theorem \ref{thm.lgv.kpaths.wt}.]The same argument as for
Proposition \ref{prop.lgv.kpaths.count} can be used here; just replace
$\operatorname*{sign}\left(  \sigma,\mathbf{p}\right)  :=\left(  -1\right)
^{\sigma}$ by $\operatorname*{sign}\left(  \sigma,\mathbf{p}\right)  =\left(
-1\right)  ^{\sigma}\cdot w\left(  \mathbf{p}\right)  $. (The only new
observation required is that when we exchange the tails of two paths in our
path tuple, the weight of the path tuple does not change. This is rather
clear: The weight of a path tuple is the product of the weights of all arcs in
all paths of the tuple\footnote{An arc will appear multiple times in the
product if it appears in multiple paths.}. When we exchange the tails of two
paths, some arcs get moved from one path to the other, but the total product
stays unchanged.)
\end{proof}

\subsubsection{Generalization to acyclic digraphs}

We can generalize Theorem \ref{thm.lgv.kpaths.wt} further. Indeed, we have
barely used anything specific to $\mathbb{Z}^{2}$ in our proofs; all we used
is that $\mathbb{Z}^{2}$ is a path-finite acyclic digraph. Thus, Theorem
\ref{thm.lgv.kpaths.wt} remains true if we replace $\mathbb{Z}^{2}$ by an
arbitrary such digraph. We thus obtain the following more general result:

\begin{theorem}
[LGV lemma, digraph weight version]\label{thm.lgv.kpaths.wt-dg}Let $K$ be a
commutative ring.

Let $D$ be a path-finite (but possibly infinite) acyclic digraph. We extend
Definition \ref{def.lgv.path-tups} to $D$ instead of $\mathbb{Z}^{2}$ (with
the obvious changes: \textquotedblleft lattice points\textquotedblright%
\ becomes \textquotedblleft vertices of $D$\textquotedblright).

For each arc $a$ of the digraph $D$, let $w\left(  a\right)  $ be an element
of $K$. We call this element $w\left(  a\right)  $ the \emph{weight} of $a$.

For each path $p$ of $D$, define the \emph{weight} $w\left(  p\right)  $ of
$p$ by%
\[
w\left(  p\right)  :=\prod_{a\text{ is an arc of }p}w\left(  a\right)  .
\]


For each path tuple $\mathbf{p}=\left(  p_{1},p_{2},\ldots,p_{k}\right)  $,
define the \emph{weight} $w\left(  \mathbf{p}\right)  $ of $\mathbf{p}$ by%
\[
w\left(  \mathbf{p}\right)  :=w\left(  p_{1}\right)  w\left(  p_{2}\right)
\cdots w\left(  p_{k}\right)  .
\]


Let $k\in\mathbb{N}$. Let $\mathbf{A}=\left(  A_{1},A_{2},\ldots,A_{k}\right)
$ and $\mathbf{B}=\left(  B_{1},B_{2},\ldots,B_{k}\right)  $ be two
$k$-vertices (i.e., two $k$-tuples of vertices of $D$). Then,%
\[
\det\left(  \left(  \sum_{p:A_{i}\rightarrow B_{j}}w\left(  p\right)  \right)
_{1\leq i\leq k,\ 1\leq j\leq k}\right)  =\sum_{\sigma\in S_{k}}\left(
-1\right)  ^{\sigma}\sum_{\substack{\mathbf{p}\text{ is a nipat}\\\text{from
}\mathbf{A}\text{ to }\sigma\left(  \mathbf{B}\right)  }}w\left(
\mathbf{p}\right)  .
\]
Here, \textquotedblleft$p:A_{i}\rightarrow B_{j}$\textquotedblright\ means
\textquotedblleft$p$ is a path from $A_{i}$ to $B_{j}$\textquotedblright.
\end{theorem}

\begin{proof}
Completely analogous to the proof of Theorem \ref{thm.lgv.kpaths.wt}.
\end{proof}

\subsubsection{The nonpermutable case}

One nice thing about the digraph $\mathbb{Z}^{2}$, however, is that in many
cases, the sum%
\[
\sum_{\sigma\in S_{k}}\left(  -1\right)  ^{\sigma}\sum_{\substack{\mathbf{p}%
\text{ is a nipat}\\\text{from }\mathbf{A}\text{ to }\sigma\left(
\mathbf{B}\right)  }}w\left(  \mathbf{p}\right)
\]
has only one nonzero addend. We have already seen this happen often in the
$k=2$ case (thanks to Proposition \ref{prop.lgv.jordan-2}). Here is the
analogous statement for general $k$:

\begin{corollary}
[LGV lemma, nonpermutable lattice weight version]\label{cor.lgv.kpaths.wt-np}%
Consider the setting of Theorem \ref{thm.lgv.kpaths.wt}, but additionally
assume that%
\begin{align}
\operatorname*{x}\left(  A_{1}\right)   &  \geq\operatorname*{x}\left(
A_{2}\right)  \geq\cdots\geq\operatorname*{x}\left(  A_{k}\right)
;\label{eq.cor.lgv.kpaths.wt-np.xA}\\
\operatorname*{y}\left(  A_{1}\right)   &  \leq\operatorname*{y}\left(
A_{2}\right)  \leq\cdots\leq\operatorname*{y}\left(  A_{k}\right)
;\label{eq.cor.lgv.kpaths.wt-np.yA}\\
\operatorname*{x}\left(  B_{1}\right)   &  \geq\operatorname*{x}\left(
B_{2}\right)  \geq\cdots\geq\operatorname*{x}\left(  B_{k}\right)
;\label{eq.cor.lgv.kpaths.wt-np.xB}\\
\operatorname*{y}\left(  B_{1}\right)   &  \leq\operatorname*{y}\left(
B_{2}\right)  \leq\cdots\leq\operatorname*{y}\left(  B_{k}\right)  .
\label{eq.cor.lgv.kpaths.wt-np.yB}%
\end{align}
Here, $\operatorname*{x}\left(  P\right)  $ and $\operatorname*{y}\left(
P\right)  $ denote the two coordinates of any point $P\in\mathbb{Z}^{2}$.

Then, there are no nipats from $\mathbf{A}$ to $\sigma\left(  \mathbf{B}%
\right)  $ when $\sigma\in S_{k}$ is not the identity permutation
$\operatorname*{id}\in S_{k}$. Therefore, the claim of Theorem
\ref{thm.lgv.kpaths.wt} simplifies to
\begin{align}
&  \det\left(  \left(  \sum_{p:A_{i}\rightarrow B_{j}}w\left(  p\right)
\right)  _{1\leq i\leq k,\ 1\leq j\leq k}\right) \nonumber\\
&  =\sum_{\substack{\mathbf{p}\text{ is a nipat}\\\text{from }\mathbf{A}\text{
to }\mathbf{B}}}w\left(  \mathbf{p}\right)  .
\label{eq.cor.lgv.kpaths.wt-np.claim}%
\end{align}

\end{corollary}

\begin{proof}
[Proof of Corollary \ref{cor.lgv.kpaths.wt-np} (sketched).]This is easy using
Proposition \ref{prop.lgv.jordan-2}. Here are the details: \medskip

\begin{fineprint}
Let $\sigma\in S_{k}$ be a permutation that is not the identity permutation
$\operatorname*{id}\in S_{k}$. Then, we don't have $\sigma\left(  1\right)
\leq\sigma\left(  2\right)  \leq\cdots\leq\sigma\left(  k\right)  $ (since
$\sigma$ is not $\operatorname*{id}$). In other words, there exists some
$i\in\left[  k-1\right]  $ such that $\sigma\left(  i\right)  >\sigma\left(
i+1\right)  $. Consider this $i$.

Now, let $\mathbf{p}$ be a nipat from $\mathbf{A}$ to $\sigma\left(
\mathbf{B}\right)  $. Write $\mathbf{p}$ in the form $\mathbf{p}=\left(
p_{1},p_{2},\ldots,p_{k}\right)  $. Thus, $p_{i}$ is a path from $A_{i}$ to
$B_{\sigma\left(  i\right)  }$, whereas $p_{i+1}$ is a path from $A_{i+1}$ to
$B_{\sigma\left(  i+1\right)  }$. Moreover, $p_{i}$ and $p_{i+1}$ have no
vertex in common (since $\mathbf{p}$ is a nipat).

The sequence $\left(  \operatorname*{x}\left(  B_{1}\right)
,\operatorname*{x}\left(  B_{2}\right)  ,\ldots,\operatorname*{x}\left(
B_{k}\right)  \right)  $ is weakly decreasing (by
(\ref{eq.cor.lgv.kpaths.wt-np.xB})). In other words, if $m$ and $n$ are two
elements of $\left[  k\right]  $ satisfying $m>n$, then $\operatorname*{x}%
\left(  B_{m}\right)  \leq\operatorname*{x}\left(  B_{n}\right)  $. Applying
this to $m=\sigma\left(  i\right)  $ and $n=\sigma\left(  i+1\right)  $, we
obtain $\operatorname*{x}\left(  B_{\sigma\left(  i\right)  }\right)
\leq\operatorname*{x}\left(  B_{\sigma\left(  i+1\right)  }\right)  $ (since
$\sigma\left(  i\right)  >\sigma\left(  i+1\right)  $). Likewise, using
(\ref{eq.cor.lgv.kpaths.wt-np.yB}), we can obtain $\operatorname*{y}\left(
B_{\sigma\left(  i\right)  }\right)  \geq\operatorname*{y}\left(
B_{\sigma\left(  i+1\right)  }\right)  $. However,
(\ref{eq.cor.lgv.kpaths.wt-np.xA}) shows that $\operatorname*{x}\left(
A_{i}\right)  \geq\operatorname*{x}\left(  A_{i+1}\right)  $. In other words,
$\operatorname*{x}\left(  A_{i+1}\right)  \leq\operatorname*{x}\left(
A_{i}\right)  $. Furthermore, (\ref{eq.cor.lgv.kpaths.wt-np.yA}) shows that
$\operatorname*{y}\left(  A_{i}\right)  \leq\operatorname*{y}\left(
A_{i+1}\right)  $. In other words, $\operatorname*{y}\left(  A_{i+1}\right)
\geq\operatorname*{y}\left(  A_{i}\right)  $.

Hence, Proposition \ref{prop.lgv.jordan-2} (applied to $A=A_{i}$,
$B=B_{\sigma\left(  i+1\right)  }$, $A^{\prime}=A_{i+1}$, $B^{\prime
}=B_{\sigma\left(  i\right)  }$, $p=p_{i}$ and $p^{\prime}=p_{i+1}$) yields
that $p_{i}$ and $p_{i+1}$ have a vertex in common. This contradicts the fact
that $p_{i}$ and $p_{i+1}$ have no vertex in common.

Forget that we fixed $\mathbf{p}$. We thus have found a contradiction for each
nipat $\mathbf{p}$ from $\mathbf{A}$ to $\sigma\left(  \mathbf{B}\right)  $.
Hence, there are no nipats from $\mathbf{A}$ to $\sigma\left(  \mathbf{B}%
\right)  $.

Forget that we fixed $\sigma$. We thus have proved that there are no nipats
from $\mathbf{A}$ to $\sigma\left(  \mathbf{B}\right)  $ when $\sigma\in
S_{k}$ is not the identity permutation $\operatorname*{id}\in S_{k}$. Hence,
if $\sigma\in S_{k}$ is not the identity permutation $\operatorname*{id}\in
S_{k}$, then%
\begin{equation}
\sum_{\substack{\mathbf{p}\text{ is a nipat}\\\text{from }\mathbf{A}\text{ to
}\sigma\left(  \mathbf{B}\right)  }}w\left(  \mathbf{p}\right)  =\left(
\text{empty sum}\right)  =0. \label{pf.cor.lgv.kpaths.wt-np.at}%
\end{equation}


Now, Theorem \ref{thm.lgv.kpaths.wt} yields%
\begin{align*}
&  \det\left(  \left(  \sum_{p:A_{i}\rightarrow B_{j}}w\left(  p\right)
\right)  _{1\leq i\leq k,\ 1\leq j\leq k}\right) \\
&  =\sum_{\sigma\in S_{k}}\left(  -1\right)  ^{\sigma}\sum
_{\substack{\mathbf{p}\text{ is a nipat}\\\text{from }\mathbf{A}\text{ to
}\sigma\left(  \mathbf{B}\right)  }}w\left(  \mathbf{p}\right) \\
&  =\underbrace{\left(  -1\right)  ^{\operatorname*{id}}}_{=1}\sum
_{\substack{\mathbf{p}\text{ is a nipat}\\\text{from }\mathbf{A}\text{ to
}\operatorname*{id}\left(  \mathbf{B}\right)  }}w\left(  \mathbf{p}\right)
+\sum_{\substack{\sigma\in S_{k};\\\sigma\neq\operatorname*{id}}}\left(
-1\right)  ^{\sigma}\underbrace{\sum_{\substack{\mathbf{p}\text{ is a
nipat}\\\text{from }\mathbf{A}\text{ to }\sigma\left(  \mathbf{B}\right)
}}w\left(  \mathbf{p}\right)  }_{\substack{=0\\\text{(by
(\ref{pf.cor.lgv.kpaths.wt-np.at}))}}}\\
&  \ \ \ \ \ \ \ \ \ \ \ \ \ \ \ \ \ \ \ \ \left(  \text{here, we have split
off the addend for }\sigma=\operatorname*{id}\text{ from the sum}\right) \\
&  =\sum_{\substack{\mathbf{p}\text{ is a nipat}\\\text{from }\mathbf{A}\text{
to }\operatorname*{id}\left(  \mathbf{B}\right)  }}w\left(  \mathbf{p}\right)
+\underbrace{\sum_{\substack{\sigma\in S_{k};\\\sigma\neq\operatorname*{id}%
}}\left(  -1\right)  ^{\sigma}0}_{=0}=\sum_{\substack{\mathbf{p}\text{ is a
nipat}\\\text{from }\mathbf{A}\text{ to }\operatorname*{id}\left(
\mathbf{B}\right)  }}w\left(  \mathbf{p}\right)  =\sum_{\substack{\mathbf{p}%
\text{ is a nipat}\\\text{from }\mathbf{A}\text{ to }\mathbf{B}}}w\left(
\mathbf{p}\right)
\end{align*}
(since $\operatorname*{id}\left(  \mathbf{B}\right)  =\mathbf{B}$). The proof
of Corollary \ref{cor.lgv.kpaths.wt-np} is now complete.
\end{fineprint}
\end{proof}

\begin{corollary}
\label{cor.lgv.binom-det-nonneg}Let $k\in\mathbb{N}$. Let $a_{1},a_{2}%
,\ldots,a_{k}$ and $b_{1},b_{2},\ldots,b_{k}$ be nonnegative integers such
that
\[
a_{1}\geq a_{2}\geq\cdots\geq a_{k}\ \ \ \ \ \ \ \ \ \ \text{and}%
\ \ \ \ \ \ \ \ \ \ b_{1}\geq b_{2}\geq\cdots\geq b_{k}.
\]
Then,%
\[
\det\left(  \left(  \dbinom{a_{i}}{b_{j}}\right)  _{1\leq i\leq k,\ 1\leq
j\leq k}\right)  \geq0.
\]

\end{corollary}

For example, if $a_{1}\geq a_{2}\geq a_{3}\geq0$ and $b_{1}\geq b_{2}\geq
b_{3}\geq0$, then%
\[
\det\left(
\begin{array}
[c]{ccc}%
\dbinom{a_{1}}{b_{1}} & \dbinom{a_{1}}{b_{2}} & \dbinom{a_{1}}{b_{3}}\\
\dbinom{a_{2}}{b_{1}} & \dbinom{a_{2}}{b_{2}} & \dbinom{a_{2}}{b_{3}}\\
\dbinom{a_{3}}{b_{1}} & \dbinom{a_{3}}{b_{2}} & \dbinom{a_{3}}{b_{3}}%
\end{array}
\right)  \geq0.
\]


\begin{proof}
[Proof of Corollary \ref{cor.lgv.binom-det-nonneg} (sketched).]Set
$K=\mathbb{Z}$, and set $w\left(  a\right)  :=1$ for each arc $a$ of
$\mathbb{Z}^{2}$. Define the lattice points%
\[
A_{i}:=\left(  0,-a_{i}\right)  \ \ \ \ \ \ \ \ \ \ \text{and}%
\ \ \ \ \ \ \ \ \ \ B_{i}:=\left(  b_{i},-b_{i}\right)
\]
for all $i\in\left[  k\right]  $. These lattice points satisfy the assumptions
of Corollary \ref{cor.lgv.kpaths.wt-np}. Hence,
(\ref{eq.cor.lgv.kpaths.wt-np.claim}) entails%
\begin{equation}
\det\left(  \left(  \sum_{p:A_{i}\rightarrow B_{j}}w\left(  p\right)  \right)
_{1\leq i\leq k,\ 1\leq j\leq k}\right)  =\sum_{\substack{\mathbf{p}\text{ is
a nipat}\\\text{from }\mathbf{A}\text{ to }\mathbf{B}}}w\left(  \mathbf{p}%
\right)  .\nonumber
\end{equation}
Since all the weights $w\left(  p\right)  $ and $w\left(  \mathbf{p}\right)  $
are $1$ in our situation, we can rewrite this as%
\begin{equation}
\det\left(  \left(  \text{\# of paths from }A_{i}\text{ to }B_{j}\right)
_{1\leq i\leq k,\ 1\leq j\leq k}\right)  =\left(  \text{\# of nipats from
}\mathbf{A}\text{ to }\mathbf{B}\right)  .\nonumber
\end{equation}


Using Proposition \ref{prop.lgv.1-paths.ct}, we can easily see that the matrix
on the left hand side of this equality is $\left(  \dbinom{a_{i}}{b_{j}%
}\right)  _{1\leq i\leq k,\ 1\leq j\leq k}$. Thus, this equality rewrites as%
\begin{equation}
\det\left(  \left(  \dbinom{a_{i}}{b_{j}}\right)  _{1\leq i\leq k,\ 1\leq
j\leq k}\right)  =\left(  \text{\# of nipats from }\mathbf{A}\text{ to
}\mathbf{B}\right)  .\nonumber
\end{equation}
Its left hand side is therefore $\geq0$ (since its right hand side is $\geq
0$). This proves Corollary \ref{cor.lgv.binom-det-nonneg}.
\end{proof}

\begin{corollary}
\label{cor.lgv.catalan-hankel-det-0}Let $k\in\mathbb{N}$. Recall the Catalan
numbers $c_{n}=\dfrac{1}{n+1}\dbinom{2n}{n}$ for all $n\in\mathbb{N}$. Then,%
\[
\det\left(  \left(  c_{i+j-2}\right)  _{1\leq i\leq k,\ 1\leq j\leq k}\right)
=\det\left(
\begin{array}
[c]{cccc}%
c_{0} & c_{1} & \cdots & c_{k-1}\\
c_{1} & c_{2} & \cdots & c_{k}\\
\vdots & \vdots & \ddots & \vdots\\
c_{k-1} & c_{k} & \cdots & c_{2k-2}%
\end{array}
\right)  =1.
\]

\end{corollary}

\begin{proof}
[Proof of Corollary \ref{cor.lgv.catalan-hankel-det-0} (sketched).]We will use
not the lattice $\mathbb{Z}^{2}$, but a different digraph. Namely, we use the
simple digraph with vertex set $\mathbb{Z}\times\mathbb{N}$ (that is, the
vertices are the lattice points that lie on the x-axis or above it) and arcs%
\[
\left(  i,j\right)  \rightarrow\left(  i+1,\ j+1\right)
\ \ \ \ \ \ \ \ \ \ \text{for all }\left(  i,j\right)  \in\mathbb{Z}%
\times\mathbb{N}%
\]
and%
\[
\left(  i,j\right)  \rightarrow\left(  i+1,\ j-1\right)
\ \ \ \ \ \ \ \ \ \ \text{for all }\left(  i,j\right)  \in\mathbb{Z}%
\times\mathbb{P},
\]
where $\mathbb{P}:=\left\{  1,2,3,\ldots\right\}  $. Here is a picture of a
small part of this digraph:%
\[%
% [tikz figure removed]
%BeginExpansion
% [tikz figure removed]%
%EndExpansion
\ \ .
\]
As we know, the Catalan number $c_{n}$ counts the paths from $\left(
0,0\right)  $ to $\left(  2n,0\right)  $ on this digraph (indeed, these are
just the Dyck paths\footnote{See Example 2 in Section \ref{sec.gf.exas} for
the definition of a Dyck path.}). Hence, $c_{n}$ also counts the paths from
$\left(  i,0\right)  $ to $\left(  2n+i,0\right)  $ whenever $i\in\mathbb{N}$
(because these are just the Dyck paths shifted by $i$ in the x-direction). It
is easy to see that this digraph is acyclic and path-finite.

Now, define two $k$-vertices $\mathbf{A}=\left(  A_{1},A_{2},\ldots
,A_{k}\right)  $ and $\mathbf{B}=\left(  B_{1},B_{2},\ldots,B_{k}\right)  $ by
setting%
\[
A_{i}:=\left(  -2\left(  i-1\right)  ,0\right)  \ \ \ \ \ \ \ \ \ \ \text{and}%
\ \ \ \ \ \ \ \ \ \ B_{i}:=\left(  2\left(  i-1\right)  ,0\right)
\]
for all $i\in\left[  k\right]  $. It is not hard to show (see Exercise
\ref{exe.lgv.catalan-hankel-det-0} \textbf{(a)}) that there is only one nipat
from $\mathbf{A}$ to $\mathbf{B}$, which is shown in the case $k=4$ on the
following picture:\footnote{In this picture, we are drawing only
\textquotedblleft half\textquotedblright\ of the grid. Indeed, the vertices
$\left(  i,j\right)  $ of our digraph $\mathbb{Z}\times\mathbb{N}$ can be
classified into \emph{even vertices} (i.e., the ones for which $i+j$ is even)
and \emph{odd vertices} (i.e., the ones for which $i+j$ is odd). Any arc
either connects two even vertices or connects two odd vertices. Hence, a path
starting at an even vertex cannot contain any odd vertex (and vice versa).
Since all our vertices $A_{1},A_{2},\ldots,A_{k}$ and $B_{1},B_{2}%
,\ldots,B_{k}$ are even, we thus don't have to bother even drawing the odd
vertices (as they have no chance to appear in any paths between our vertices).
As a consequence, we are drawing only the grid lines containing the even
vertices.}%
\[%
% [tikz figure removed]
%BeginExpansion
% [tikz figure removed]%
%EndExpansion
\]
(the point $A_{1}$ coincides with $B_{1}$, and the path from $A_{1}$ to
$B_{1}$ is invisible, since it has no arcs). Moreover, it can be shown (see
Exercise \ref{exe.lgv.catalan-hankel-det-0} \textbf{(b)}) that there are no
nipats from $\mathbf{A}$ to $\sigma\left(  \mathbf{B}\right)  $ when
$\sigma\in S_{k}$ is not the identity permutation $\operatorname*{id}\in
S_{k}$. (This is analogous to Corollary \ref{cor.lgv.kpaths.wt-np}.) Hence, if
we set $K=\mathbb{Z}$ and $w\left(  a\right)  =1$ for each arc $a$ of our
digraph, then (\ref{eq.cor.lgv.kpaths.wt-np.claim}) entails%
\begin{equation}
\det\left(  \left(  \sum_{p:A_{i}\rightarrow B_{j}}w\left(  p\right)  \right)
_{1\leq i\leq k,\ 1\leq j\leq k}\right)  =\sum_{\substack{\mathbf{p}\text{ is
a nipat}\\\text{from }\mathbf{A}\text{ to }\mathbf{B}}}w\left(  \mathbf{p}%
\right) \nonumber
\end{equation}
(by the same reasoning as in the proof of Corollary \ref{cor.lgv.kpaths.wt-np}%
). The right hand side of this equality is $1$ (since there is only one nipat
from $\mathbf{A}$ to $\mathbf{B}$), while the matrix on the left hand side is
easily seen to be $\left(  c_{i+j-2}\right)  _{1\leq i\leq k,\ 1\leq j\leq k}$
(since the \# of paths from $A_{i}$ to $B_{j}$ is the Catalan number
$c_{i+j-2}$). This yields the claim of Corollary
\ref{cor.lgv.catalan-hankel-det-0}. The details are LTTR.
\end{proof}

The LGV lemma in all its variants is one major place in which combinatorial
questions reduce to the computation of determinants. Other such places are the
\emph{matrix-tree theorem} (see, e.g., \cite[\S 4]{Zeilbe}, \cite[\S 3.17]%
{Loehr-BC}, \cite[Theorems 9.8 and 10.4]{Stanle18}, \cite[\S 5.14]{23s}) and
the enumeration of perfect matchings or domino tilings (see, e.g.,
\cite{Stucky15}, \cite[\S 12.12--\S 12.13]{Loehr-BC}, \cite[\S 10.1]%
{Aigner07}). Soon, we will also encounter determinants in the study of
symmetric functions.

See also \cite{BruRys91} for a selection of other intersections between
combinatorics and linear algebra.

\begin{noncompile}
TODO: Cover the matrix-tree theorem (OH8).
\end{noncompile}

\section{\label{chap.sf}Symmetric functions}

This final chapter is devoted to the theory of \emph{symmetric functions}.
Specifically, we will restrict ourselves to \emph{symmetric polynomials} (the
\textquotedblleft functions\textquotedblright\ part is a technical tweak that
makes the theory neater but we won't have time to introduce). Serious
treatments of the subject can be found in \cite{Wildon20}, \cite[Chapters
10--11]{Loehr-BC}, \cite{Egge19}, \cite{MenRem15}, \cite{Macdon95},
\cite[Chapter 8]{Aigner07}, \cite[Chapter 7]{Stanley-EC2}, \cite[Chapter
7]{Sagan19}, \cite[Chapter 4]{Sagan01}, \cite{Krishn86}, \cite[Chapters
14--19]{FoaHan04}, \cite{Savage22}, \cite{SmiTut24}, \cite[Chapter 2]{GriRei}
and \cite{LLPT95}.

We begin with some oversimplified historical motivation.

Symmetric polynomials first(?) appeared in the study of roots of polynomials.
Consider a monic univariate polynomial%
\[
f=x^{n}+a_{1}x^{n-1}+a_{2}x^{n-2}+\cdots+a_{n}x^{0}\in\mathbb{C}\left[
x\right]
\]
(note the nonstandard labeling of coefficients). Let $r_{1},r_{2},\ldots
,r_{n}\in\mathbb{C}$ be the roots of this polynomial (listed with
multiplicities). Then, Fran\c{c}ois Vi\`{e}te (aka Franciscus Vieta) noticed
that%
\[
f=\left(  x-r_{1}\right)  \left(  x-r_{2}\right)  \cdots\left(  x-r_{n}%
\right)  ,
\]
so that (by comparing coefficients) we see that%
\begin{align*}
r_{1}+r_{2}+\cdots+r_{n}  &  =-a_{1};\\
\sum_{i<j}r_{i}r_{j}  &  =a_{2};\\
\sum_{i<j<k}r_{i}r_{j}r_{k}  &  =-a_{3};\\
&  \ldots;\\
r_{1}r_{2}\cdots r_{n}  &  =\left(  -1\right)  ^{n}a_{n}.
\end{align*}
These equalities are now known as \emph{Viete's formulas}. They allow
computing certain expressions in the $r_{i}$'s without having to compute the
$r_{i}$'s themselves. For instance, we can compute $r_{1}^{2}+r_{2}^{2}%
+\cdots+r_{n}^{2}$ (that is, the sum of the squares of all roots of $f$) by
observing that%
\[
\left(  r_{1}+r_{2}+\cdots+r_{n}\right)  ^{2}=\left(  r_{1}^{2}+r_{2}%
^{2}+\cdots+r_{n}^{2}\right)  +2\sum_{i<j}r_{i}r_{j},
\]
so that%
\[
r_{1}^{2}+r_{2}^{2}+\cdots+r_{n}^{2}=\left(  \underbrace{r_{1}+r_{2}%
+\cdots+r_{n}}_{=-a_{1}}\right)  ^{2}-2\underbrace{\sum_{i<j}r_{i}r_{j}%
}_{=a_{2}}=a_{1}^{2}-2a_{2}.
\]
This shows, among other things, that $r_{1}^{2}+r_{2}^{2}+\cdots+r_{n}^{2}$ is
an integer if the coefficients of $f$ are integers. Newton and others found
similar formulas for $r_{1}^{3}+r_{2}^{3}+\cdots+r_{n}^{3}$ and other such
polynomials. (These formulas are now known as the \emph{Newton--Girard
identities} -- see Theorem \ref{thm.sf.NG} below.) Gauss extended this to
arbitrary symmetric polynomials in $r_{1},r_{2},\ldots,r_{n}$ (by
algorithmically expressing them as polynomials in $a_{1},a_{2},\ldots,a_{n}$),
and used it in one of his proofs of the Fundamental Theorem of Algebra
\cite{Gauss16}; Galois used this to build what is now known as Galois theory
(even though modern treatments of Galois theory often avoid symmetric
polynomials); some harbingers of this can be seen in Cardano's solution of the
cubic equation. See \cite{Tignol16} and \cite{Armstrong-561fa18sp19} for the
real history.

Here is a simple modern application of the same ideas: Let $A\in
\mathbb{C}^{n\times n}$ be a matrix with eigenvalues $\lambda_{1},\lambda
_{2},\ldots,\lambda_{n}$ (listed with algebraic multiplicities). Let
$f\in\mathbb{C}\left[  x\right]  $ be a univariate polynomial. The
\emph{spectral mapping theorem} says that the eigenvalues of the matrix
$f\left[  A\right]  $ are $f\left[  \lambda_{1}\right]  ,f\left[  \lambda
_{2}\right]  ,\ldots,f\left[  \lambda_{n}\right]  $ (here, I am using the
notation $f\left[  a\right]  $ for the value of $f$ at some element $a$; this
is usually written $f\left(  a\right)  $). Thus, the characteristic polynomial
of $f\left[  A\right]  $ is
\begin{align*}
\chi_{f\left[  A\right]  }  &  =\left(  x-f\left[  \lambda_{1}\right]
\right)  \left(  x-f\left[  \lambda_{2}\right]  \right)  \cdots\left(
x-f\left[  \lambda_{n}\right]  \right) \\
&  =x^{n}-\left(  f\left[  \lambda_{1}\right]  +f\left[  \lambda_{2}\right]
+\cdots+f\left[  \lambda_{n}\right]  \right)  x^{n-1}+\left(  \sum
_{i<j}f\left[  \lambda_{i}\right]  f\left[  \lambda_{j}\right]  \right)
x^{n-2}\pm\cdots.
\end{align*}
Hence, all coefficients of $\chi_{f\left[  A\right]  }$ are symmetric
polynomials in the $\lambda_{i}$s that depend only on $f$ (not on $A$). In
particular, this shows that $\chi_{f\left[  A\right]  }$ is uniquely
determined by $f$ and $\chi_{A}$. But can you compute $\chi_{f\left[
A\right]  }$ exactly in terms of $f$ and $\chi_{A}$ without computing the
roots $\lambda_{1},\lambda_{2},\ldots,\lambda_{n}$ ? Yes, if you know how to
express \textbf{any} symmetric polynomial in the $\lambda_{i}$s in terms of
the coefficients of $\chi_{A}$. This is the same problem that Gauss solved
with his algorithm for expressing an arbitrary symmetric polynomial in the
roots of a polynomial in terms of the coefficients of the polynomial.
Incidentally, this algorithm also becomes helpful when one tries to generalize
the spectral mapping theorem to matrices over arbitrary commutative rings.
Here, eigenvalues don't always exist (let alone $n$ of them), so it becomes
necessary to restate the theorem in a language that does not rely on them.
Again, symmetric polynomials provide the way to do this.
